\chapter{Singularities and the Limits of Their Dilations}

\section*{1. Classifying maximal solutions}

\begin{remark}[Singularity time and maximal solutions]
\label{ch08-rem-singularity-time-maximal-solutions}
\source{chow-knopf-ricci-flow:page-0246}
\dcref{Remark 8.0}
\group{classifying-maximal-solutions}

If \((\mathcal M^n,g(t))\) is a Ricci flow solution defined on a maximal time interval \([0,T)\) with \(T\in(0,\infty]\), then \(T\) is the singularity time. The chapter distinguishes finite-time singularities (\(T<\infty\)) from immortal solutions (\(T=\infty\)).
\end{remark}

\begin{definition}[Type I, IIa, IIb, and III singularities]
\label{ch08-def-singularity-types}
\source{chow-knopf-ricci-flow:page-0247}
\dcref{Definition 8.1}
\group{classifying-maximal-solutions}

Let \((\mathcal M^n,g(t))\) be a solution of the Ricci flow on a maximal time interval \( [0,T)\), where \(T\in(0,\infty]\).
\begin{itemize}
\item It encounters a \textbf{Type I singularity} if \(T<\infty\) and
\[
\sup_{\mathcal M^n\times[0,T)} |\Rm(\cdot,t)|(T-t)<\infty.
\]
\item It encounters a \textbf{Type IIa singularity} if \(T<\infty\) and
\[
\sup_{\mathcal M^n\times[0,T)} |\Rm(\cdot,t)|(T-t)=\infty.
\]
\item It encounters a \textbf{Type IIb singularity} if \(T=\infty\) and
\[
\sup_{\mathcal M^n\times[0,\infty)} |\Rm(\cdot,t)|\,t=\infty.
\]
\item It encounters a \textbf{Type III singularity} if \(T=\infty\) and
\[
\sup_{\mathcal M^n\times[0,\infty)} |\Rm(\cdot,t)|\,t<\infty.
\]
\end{itemize}
\end{definition}

\section*{2. Singularity models}

\begin{remark}[Singularity models and limit solutions]
\label{ch08-rem-singularity-models}
\source{chow-knopf-ricci-flow:page-0248}
\dcref{Remark 8.2}
\group{singularity-models}

Singularity models are complete nonflat limit solutions obtained as limits of dilations about singularities. They exist on infinite time intervals, and each singularity type has an associated model class.
\end{remark}

\begin{definition}[Ancient, eternal, and immortal singularity models]
\label{ch08-def-model-types}
\source{chow-knopf-ricci-flow:page-0248, chow-knopf-ricci-flow:page-0249}
\dcref{Definition 8.3}
\group{singularity-models}

Let \((\mathcal M_\infty^n,g_\infty(t))\) be a limit solution of the Ricci flow.
\begin{itemize}
\item It is an \textbf{ancient Type I singularity model} if it exists on \((-\infty,\omega)\) containing \(t=0\) and satisfies
\[
\sup_{\mathcal M_\infty^n\times(-\infty,0]} |\Rm_\infty(\cdot,t)|\,|t|<\infty.
\]
\item It is an \textbf{ancient Type II singularity model} if it exists on \((-\infty,\omega)\) containing \(t=0\) and satisfies
\[
\sup_{\mathcal M_\infty^n\times(-\infty,0]} |\Rm_\infty(\cdot,t)|\,|t|=\infty.
\]
\item It is an \textbf{eternal Type II singularity model} if it exists for all \(t\in(-\infty,\infty)\) and satisfies
\[
\sup_{\mathcal M_\infty^n\times(-\infty,\infty)} |\Rm_\infty(\cdot,t)|<\infty.
\]
\item It is an \textbf{immortal Type III singularity model} if it exists on \((-\alpha,\infty)\) containing \(t=0\) and satisfies
\[
\sup_{\mathcal M_\infty^n\times[0,\infty)} |\Rm_\infty(\cdot,t)|\,t<\infty.
\]
\end{itemize}
\end{definition}

\begin{remark}[Point picking for singularity models]
\label{ch08-rem-point-picking}
\source{chow-knopf-ricci-flow:page-0249}
\dcref{Remark 8.4}
\group{singularity-models}

By choosing the dilating sequence carefully, one can arrange that a singularity model satisfies
\[
\sup_{x\in \mathcal M^n} |\Rm_\infty(x,0)|=|\Rm_\infty(y,0)|
\]
for some \(y\in\mathcal M^n\). If the curvature operator is nonnegative, one can instead arrange
\[
\sup_{x\in \mathcal M^n} R_\infty(x,0)=R_\infty(y,0).
\]
\end{remark}

\begin{remark}[Dimension three and scalar curvature]
\label{ch08-rem-dimension-three-scalar-curvature}
\source{chow-knopf-ricci-flow:page-0249}
\dcref{Remark 8.5}
\group{singularity-models}

By Theorem 9.4, any limit of a finite-time singularity in dimension three has nonnegative curvature operator. Consequently, in dimension three one may always satisfy the scalar-curvature version of the point-picking condition.
\end{remark}

\begin{remark}[Ancient versus eternal Type II limits]
\label{ch08-rem-ancient-eternal-type-ii}
\source{chow-knopf-ricci-flow:page-0250}
\dcref{Remark 8.6}
\group{singularity-models}

The distinction between ancient and eternal Type II models matters because any limit of a Type II singularity is noncompact. To study geometry at spatial infinity, one may take a second limit by dimension reduction, which produces ancient solutions of Type I or Type II.
\end{remark}

\begin{example}[The Rosenau solution]
\label{ch08-ex-rosenau-solution}
\source{chow-knopf-ricci-flow:page-0250}
\dcref{Example 8.7}
\group{singularity-models}

The Rosenau solution is an example of a Type II ancient solution that is not eternal, but whose backward limit is an eternal solution, namely the cigar soliton.
\end{example}

\begin{remark}[Finite-time singularities in the Ricci flow program]
\label{ch08-rem-finite-time-singularities}
\source{chow-knopf-ricci-flow:page-0250}
\dcref{Remark 8.8}
\group{singularity-models}

Finite-time singularities, namely Type I and Type IIa, are the singularities most relevant to the Ricci flow program for closed \(3\)-manifolds.
\end{remark}

\begin{example}[Neckpinch and degenerate neckpinch]
\label{ch08-ex-neckpinch-degenerate-neckpinch}
\source{chow-knopf-ricci-flow:page-0250}
\dcref{Example 8.9}
\group{singularity-models}

In dimension \(3\), the neckpinch and the conjectured degenerate neckpinch are the canonical examples of Type I and Type IIa singularities, respectively.
\end{example}

\section*{8. Dilations of singularities}

\begin{definition}[Globally curvature essential sequences]
\label{ch08-def-globally-curvature-essential}
\source{chow-knopf-ricci-flow:page-0251}
\dcref{Definition 8.10}
\group{dilations-of-singularities}

A sequence \((x_i,t_i)\) is \textbf{globally curvature essential} if \(t_i\nearrow T\le\infty\) and there exist \(C\ge1\) and radii \(r_i\in(0,\sqrt{t_i})\) such that
\[
\sup\Bigl\{ |\Rm(x,t)| : x\in \overline B_{g(t)}(x_i,r_i),\ t\in[t_i-r_i^2,t_i]\Bigr\}
\le C |\Rm(x_i,t_i)|
\]
for all \(i\), and
\[
\lim_{i\to\infty} r_i^2 |\Rm(x_i,t_i)|=\infty.
\]
\end{definition}

\begin{definition}[Locally curvature essential sequences]
\label{ch08-def-locally-curvature-essential}
\source{chow-knopf-ricci-flow:page-0251}
\dcref{Definition 8.11}
\group{dilations-of-singularities}

A sequence \((x_i,t_i)\) is \textbf{locally curvature essential} if \(t_i\nearrow T\le\infty\) and there exist \(C<\infty\) and radii \(r_i\in(0,\sqrt{t_i})\) such that
\[
\sup\Bigl\{ |\Rm(x,t)| : x\in \overline B_{g(t)}(x_i,r_i),\ t\in[t_i-r_i^2,t_i]\Bigr\}
\le C r_i^{-2}
\]
for all \(i\).
\end{definition}

\begin{remark}[Globally essential sequences and limits]
\label{ch08-rem-globally-essential-sequences}
\source{chow-knopf-ricci-flow:page-0251}
\dcref{Remark 8.12}
\group{dilations-of-singularities}

In some arguments one uses the special case
\[
\sup\Bigl\{ |\Rm(x,t)| : x\in \mathcal M^n,\ t\in[0,t_i]\Bigr\}\le C|\Rm(x_i,t_i)|
\]
of global curvature essentiality.
\end{remark}

\begin{remark}[Consequences of the essentiality hypotheses]
\label{ch08-rem-essentiality-consequences}
\source{chow-knopf-ricci-flow:page-0251}
\dcref{Remark 8.13}
\group{dilations-of-singularities}

Any smooth limit of a globally curvature essential sequence formed by the parabolic dilations below has curvature bounded by a function of time, and the limit is complete. A smooth limit of a locally curvature essential sequence has bounded curvature on a unit metric ball for times \(-1\le t\le0\).
\end{remark}

\begin{definition}[Parabolic dilations]
\label{ch08-def-parabolic-dilations}
\source{chow-knopf-ricci-flow:page-0252}
\dcref{Equation 8.8}
\group{dilations-of-singularities}

Given a sequence \((x_i,t_i)\) with \(t_i\nearrow T\in(0,\infty]\), define the parabolic dilations
\[
g_i(t)\doteq |\Rm(x_i,t_i)|\cdot g\!\left(t_i+\frac{t}{|\Rm(x_i,t_i)|}\right),
\]
which exist for
\[
-t_i|\Rm(x_i,t_i)|\le t < (T-t_i)|\Rm(x_i,t_i)|.
\]
\end{definition}

\begin{remark}[Normalization of curvature under dilation]
\label{ch08-rem-dilation-normalization}
\source{chow-knopf-ricci-flow:page-0252}
\dcref{Remark 8.14}
\group{dilations-of-singularities}

The dilation is chosen so that the curvature of \(g_i\) satisfies \(|\Rm[g_i](x_i,0)|_{g_i}=1\). In dimension \(3\), one may replace \(|\Rm|\) by \(R\) in the dilation discussion.
\end{remark}

\begin{definition}[Global injectivity radius estimate]
\label{ch08-def-global-injectivity-radius-estimate}
\source{chow-knopf-ricci-flow:page-0252}
\dcref{Definition 8.15}
\group{dilations-of-singularities}

A solution \((\mathcal M^n,g(t))\) on \([0,T)\) satisfies a \textbf{global injectivity radius estimate on the scale of its maximum curvature} if there exists \(c>0\) such that
\[
\inj(x,t)^2\ge \frac{c}{\sup_{\mathcal M^n}|\Rm(\cdot,t)|}
\]
for all \((x,t)\in\mathcal M^n\times[0,T)\).
\end{definition}

\begin{definition}[Kappa-noncollapsedness]
\label{ch08-def-kappa-noncollapsed}
\source{chow-knopf-ricci-flow:page-0253}
\dcref{Definition 8.16}
\group{dilations-of-singularities}

A solution \((\mathcal M^n,g(t))\) on \([0,T)\) is \(\kappa\)-noncollapsed on the scale \(\rho\) if every ball \(B_{g(t)}(x_0,r)\) with \(r<\rho\) and curvature bounded by \(r^{-2}\) on the corresponding parabolic neighborhood satisfies
\[
\Vol[B_{g(t)}(x_0,r)]\ge \kappa r^n.
\]
\end{definition}

\begin{theorem}[Perelman's injectivity estimate for finite-time singularities]
\label{ch08-thm-perelman-injectivity-estimate}
\source{chow-knopf-ricci-flow:page-0253}
\dcref{Theorem 8.17}
\group{dilations-of-singularities}

If \((\mathcal M^n,g(t))\) encounters a singularity at some finite time \(T<\infty\), then there exists \(c>0\) and a subsequence \((x_i,t_i)\) such that
\[
\inj(x_i,t_i)\ge \frac{c}{\sqrt{\max_{\mathcal M^n}|\Rm(\cdot,t_i)|}}.
\]
\end{theorem}

\begin{proof}
\source{chow-knopf-ricci-flow:page-0253}
\sketch The source states this injectivity-radius estimate as a theorem and indicates that a detailed proof will be given in later work. We record the theorem here as an input to the chapter's dilation arguments.
\end{proof}

\begin{remark}[Hamilton's earlier estimate]
\label{ch08-rem-hamilton-injectivity}
\source{chow-knopf-ricci-flow:page-0253}
\dcref{Remark 8.18}
\group{dilations-of-singularities}

For Type I singularities, a global injectivity radius estimate on the scale of the maximum curvature was originally obtained from Hamilton's isoperimetric estimate.
\end{remark}

\section*{4. Dilations of finite-time singularities}

\begin{lemma}[Lower blow-up rate for finite-time singularities]
\label{ch08-lem-finite-time-blowup-rate}
\source{chow-knopf-ricci-flow:page-0253, chow-knopf-ricci-flow:page-0254}
\dcref{Lemma 8.19}
\group{dilations-of-finite-time-singularities}

Let \(\mathcal M^n\) be compact, and let \(g(t)\) be a maximal Ricci flow on \([0,T)\) with \(T<\infty\). Then there exists \(c_0>0\), depending only on \(n\), such that
\[
\sup_{x\in\mathcal M^n}|\Rm(x,t)|\ge \frac{c_0}{T-t}
\]
for all \(t\in[0,T)\).
\end{lemma}

\begin{proof}
\source{chow-knopf-ricci-flow:page-0254}
\sketch Let \(K(t)=\sup_{x\in\mathcal M^n}|\Rm(x,t)|^2\). By the curvature evolution inequality from Chapter 7,
\[
\partial_t |\Rm|^2\le \Delta |\Rm|^2+C|\Rm|^3,
\]
and the maximum principle gives \(K'(t)\le C K(t)^{3/2}\). Equivalently,
\[
\frac{d}{dt}K(t)^{-1/2}\ge -\frac C2.
\]
Since \(K(t)\to\infty\) as \(t\nearrow T\), one has \(\liminf_{\tau\to T}K(\tau)^{-1/2}=0\). Integrating from \(t\) to \(\tau\) and sending \(\tau\nearrow T\) yields
\[
K(t)^{-1/2}\le \frac C2(T-t),
\]
which implies the stated lower bound after renaming constants.
\end{proof}

\begin{proposition}[Type I limits of Type I singularities]
\label{ch08-prop-type-i-limits}
\source{chow-knopf-ricci-flow:page-0255, chow-knopf-ricci-flow:page-0256}
\dcref{Proposition 8.20}
\group{dilations-of-finite-time-singularities}

If \((\mathcal M^n,g(t))\) has a Type I singularity at \(T\in(0,\infty)\) and \((x_i,t_i)\) is a globally curvature essential sequence, then any pointed limit of the parabolic dilations is an ancient Type I singularity model. Moreover, that limit itself exhibits a Type I singularity at some finite time \(\omega<\infty\).
\end{proposition}

\begin{proof}
\source{chow-knopf-ricci-flow:page-0255, chow-knopf-ricci-flow:page-0256}
\sketch Let
\[
S=\sup_{\mathcal M^n\times[0,T)}|\Rm(x,t)|(T-t)<\infty.
\]
The Type I hypothesis and Lemma~\dcref{Lemma 8.19} give constants \(c_0<C\) with
\[
\frac{c_0}{T-t}\le \sup_{\mathcal M^n}|\Rm(\cdot,t)|\le \frac{C}{T-t}.
\]
By global curvature essentiality, \(|\Rm(x_i,t_i)|(T-t_i)\) stays uniformly comparable to \(\sup_{\mathcal M^n}|\Rm(\cdot,t_i)|(T-t_i)\), so the dilation factors
\[
\omega_i=(T-t_i)|\Rm(x_i,t_i)|
\]
stay in a compact subset of \((0,\infty)\). After passing to a subsequence, \(\omega_i\to\omega\in[c,C]\). The same estimates show that \(\alpha_i=t_i|\Rm(x_i,t_i)|\to\infty\). The parabolic dilations therefore exist on longer and longer intervals \([-\alpha_i,\omega_i)\), and the curvature bounds obtained from the Type I hypothesis yield a pointed limit defined on \((-\infty,\omega)\) with the same Type I asymptotics at time \(\omega\).
\end{proof}

\begin{remark}[Shrinking spherical space forms in dimension three]
\label{ch08-rem-shrinking-spherical-space-forms}
\source{chow-knopf-ricci-flow:page-0256}
\dcref{Remark 8.21}
\group{dilations-of-finite-time-singularities}

Any Type I singularity starting from a compact \(3\)-manifold of positive Ricci curvature yields a Type I limit that is a shrinking spherical space form.
\end{remark}

\begin{remark}[Relative curvature scale]
\label{ch08-rem-relative-curvature-scale}
\source{chow-knopf-ricci-flow:page-0256}
\dcref{Equation 8.13}
\group{dilations-of-finite-time-singularities}

The relative curvature scale of a point \(x\) at time \(t\) is the ratio
\[
\frac{|\Rm(x,t)|}{\sup_{\mathcal M^n}|\Rm(\cdot,t)|}.
\]
\end{remark}

\begin{definition}[The Type I limsup scale]
\label{ch08-def-type-i-limsup-scale}
\source{chow-knopf-ricci-flow:page-0256}
\dcref{Equation 8.14}
\group{dilations-of-finite-time-singularities}

Define
\[
\omega\coloneqq \limsup_{t\uparrow T}|\Rm(\cdot,t)|(T-t).
\]
\end{definition}

\begin{definition}[Type I rescaling factors]
\label{ch08-def-type-i-rescaling-factors}
\source{chow-knopf-ricci-flow:page-0256, chow-knopf-ricci-flow:page-0257}
\dcref{Equation 8.15}
\group{dilations-of-finite-time-singularities}

Choose \((x_i,t_i)\) with \(t_i\uparrow T\) so that
\[
|\Rm(x_i,t_i)|(T-t_i)=:\omega_i\to\omega.
\]
\end{definition}

\begin{remark}[Type I curvature bound after rescaling]
\label{ch08-rem-type-i-curvature-bound-after-rescaling}
\source{chow-knopf-ricci-flow:page-0256, chow-knopf-ricci-flow:page-0257}
\dcref{Equation 8.16}
\group{dilations-of-finite-time-singularities}

For each \(\epsilon>0\), once \(t_i\) is sufficiently close to \(T\) one has
\[
|\Rm_i(x,t)|\le \frac{\omega+\epsilon}{\omega_i-t}
\]
for all \(x\in\mathcal M^n\) on the dilated time interval where the dilation is defined.
\end{remark}

\section*{5. Dilations of infinite-time singularities}

\subsection*{5.1. Type II limits of Type IIb singularities}

\begin{remark}[Type IIb singularity condition]
\label{ch08-rem-type-iib-condition}
\source{chow-knopf-ricci-flow:page-0259}
\dcref{Equation 8.21}
\group{dilations-of-infinite-time-singularities}

A solution forms a Type IIb singularity at \(T=\infty\) if
\[
\sup_{\mathcal M^n\times[0,\infty)} t\,|\Rm(x,t)|=\infty.
\]
\end{remark}

\begin{remark}[Refined Type IIb rescaling]
\label{ch08-rem-refined-type-iib-rescaling}
\source{chow-knopf-ricci-flow:page-0260}
\dcref{Equations 8.22, 8.23}
\group{dilations-of-infinite-time-singularities}

Choosing \((x_i,t_i)\) to nearly maximize
\[
\frac{t_i(T_i-t_i)|\Rm(x_i,t_i)|}{\sup_{\mathcal M^n\times[0,T_i]}|t|(T_i-t)|\Rm(x,t)|}
\]
ensures that both \(\alpha_i=t_i|\Rm(x_i,t_i)|\) and \(\omega_i=(T_i-t_i)|\Rm(x_i,t_i)|\) tend to \(\infty\), and the dilations satisfy a uniform curvature bound on larger and larger time intervals.
\end{remark}

\subsection*{5.2. Type III limits of Type III singularities}

\begin{definition}[The Type III limsup scale]
\label{ch08-def-type-iii-limsup-scale}
\source{chow-knopf-ricci-flow:page-0260}
\dcref{Equation 8.23}
\group{dilations-of-infinite-time-singularities}

For a Type III solution on \([0,\infty)\), define
\[
\alpha\coloneqq \limsup_{t\to\infty}\left(t\sup_{\mathcal M^n}|\Rm(\cdot,t)|\right)\in[0,\infty).
\]
\end{definition}

\begin{corollary}[Length under Ricci flow]
\label{ch08-cor-length-under-ricci-flow}
\source{chow-knopf-ricci-flow:page-0261}
\dcref{Corollary 8.24}
\group{dilations-of-infinite-time-singularities}

If \(\gamma_t\) is a time-dependent family of curves with fixed endpoints in a Ricci flow solution, and if each \(\gamma_t\) is either independent of time or geodesic with respect to \(g(t)\), then
\[
\frac{d}{dt}L_t(\gamma_t)=-\int_{\gamma_t}\Rc(S,S)\,ds.
\]
\end{corollary}

\begin{proof}
\source{chow-knopf-ricci-flow:page-0261}
\sketch The general length variation formula gives
\[
\frac{d}{dt}L_t(\gamma_t)
= -\int_{\gamma_t}\Rc(S,S)\,ds
- \int_{\gamma_t}\langle \nabla_S S,V\rangle\,ds.
\]
If \(\gamma_t\) is time-independent or geodesic with respect to \(g(t)\), then the second integral vanishes. The displayed formula follows.
\end{proof}

\begin{proposition}[Non-nilmanifold Type III lower bound]
\label{ch08-prop-type-iii-lower-bound}
\source{chow-knopf-ricci-flow:page-0261, chow-knopf-ricci-flow:page-0262}
\dcref{Proposition 8.25}
\group{dilations-of-infinite-time-singularities}

There exists \(\varepsilon>0\), depending only on \(n\), such that if \((\mathcal M^n,g(t))\) is a Ricci flow on \([0,\infty)\) and \(\alpha\le\varepsilon\), then \(\mathcal M^n\) is a nilmanifold. Equivalently, if \(\mathcal M^n\) is not a nilmanifold, then \(\alpha\ge\varepsilon\).
\end{proposition}

\begin{proof}
\source{chow-knopf-ricci-flow:page-0261, chow-knopf-ricci-flow:page-0262}
\sketch Assume \(\alpha\) is sufficiently small. Then for large \(t\), the curvature satisfies
\[
\sup_{\mathcal M^n}|\Rc(\cdot,t)|\le \frac{C\varepsilon}{t},
\]
so the length \(L(t)\) of any fixed path satisfies
\[
\frac{dL}{dt}\le \frac{C\varepsilon}{t}L.
\]
Integrating gives \(L(t)\le C' t^{C\varepsilon}\). Choosing \(\varepsilon<1/(2C)\) yields
\[
\diam(\mathcal M^n,g(t))\le Ct^{1/2-\delta}
\]
for some \(\delta>0\). Hence
\[
\sup_{\mathcal M^n}|\Rm(\cdot,t)|\,\diam(\mathcal M^n,g(t))^2\to0,
\]
and the nilmanifold conclusion follows from the geometric classification invoked in the source.
\end{proof}

\begin{remark}[Type III limits]
\label{ch08-rem-type-iii-limits}
\source{chow-knopf-ricci-flow:page-0258, chow-knopf-ricci-flow:page-0259}
\dcref{Remark 8.22}
\group{dilations-of-infinite-time-singularities}

In dimension \(3\), the estimate \(C|\Rm|-C'\le R\le C_3|\Rm|\) allows one to replace \(|\Rm|\) by scalar curvature when studying Type III limits. This is useful because the scalar curvature of a pointed Type III limit can be arranged to attain its maximum at the base point.
\end{remark}

\begin{example}[Bryant soliton]
\label{ch08-ex-bryant-soliton}
\source{chow-knopf-ricci-flow:page-0259}
\dcref{Example 8.23}
\group{dilations-of-infinite-time-singularities}

The Bryant soliton on \(\mathbb R^3\) is an example of a rotationally symmetric Type II limit.
\end{example}

\begin{remark}[Limits obtained by a third dilation]
\label{ch08-rem-third-dilation}
\source{chow-knopf-ricci-flow:page-0263, chow-knopf-ricci-flow:page-0264}
\dcref{Notes 8.26}
\group{dilations-of-infinite-time-singularities}

Starting from a Type II ancient solution, one may choose points and times \(t_i\searrow -\infty\) and dilate again to obtain an eternal pointed limit that satisfies
\[
\sup |\Rm_{2\infty}|\le 1 = |\Rm_{2\infty}(x_{2\infty},0)|.
\]
\end{remark}

\begin{remark}[Chapter reference]
\label{ch08-rem-notes-reference}
\dcref{Notes 8.27}
\group{notes-and-commentary}
\source{chow-knopf-ricci-flow:page-0264}
The main reference for this chapter is Section 16 of the cited source. The limit solutions are obtained using the Compactness Theorem from Chapter 7.
\end{remark}
